\section{Objetivos}

El objetivo principal de este trabajo es diseñar e implementar una arquitectura de
código abierto, modular y reproducible para la producción remota de vídeo en tiempo real,
que resulte viable tanto en entornos de desarrollo como en producción profesional.
Partiendo del mezclador Voctomix como base, se busca construir un sistema completo que
democratice el acceso a herramientas de producción audiovisual de calidad profesional,
eliminando la dependencia de hardware propietario de alto coste y permitiendo su despliegue
mediante tecnologías de contenedores estándar.

Para alcanzar este objetivo principal se definen los siguientes objetivos específicos:

\begin{enumerate}
  \item Extender las capacidades del mezclador Voctomix con funcionalidades propias de
        la producción profesional: un sistema de rotulación dinámica con gráficos y
        textos en tiempo real, un gestor de estados de emisión (\textit{stream blanker})
        y un mecanismo de sincronización automática de audio (\textit{Audio Follows Video}).

  \item Implementar un servicio de telemetría que exporte el estado completo del sistema
        en formato JSON y lo publique en un broker de mensajería AMQP, posibilitando la
        integración con sistemas externos de monitorización sin modificar el núcleo de
        la aplicación.

  \item Contenerizar el sistema completo mediante Docker y Docker Compose, garantizando
        la reproducibilidad del entorno de ejecución y permitiendo el despliegue íntegro
        con un único comando.

  \item Extender el despliegue a una plataforma de orquestación Kubernetes, validando
        la viabilidad del sistema en infraestructuras multi-nodo y sentando las bases
        para su escalabilidad en entornos cloud.

  \item Validar el sistema en múltiples escenarios de despliegue, monopuesto nativo,
        cliente-servidor en red local, Docker Compose y Kubernetes, documentando el
        rendimiento y la cobertura funcional en cada uno.

  \item Demostrar la viabilidad operativa del sistema en un caso de uso real dentro del
        proyecto europeo CyberNEMO de la Universidad Politécnica de Madrid, validando
        su uso más allá del entorno de laboratorio en condiciones de producción reales.
\end{enumerate}

\begin{table}[H]
  \centering
  \caption{Relación entre objetivos específicos y secciones de la memoria.}
  \label{tab:objetivos_secciones}
  \begin{tabular}{cll}
    \hline
    \textbf{Obj.} & \textbf{Descripción} & \textbf{Sección} \\
    \hline
    OE-1 & Rotulación, stream blanker y AFV  & \S\ref{sec:overlays}, \S\ref{sec:streamblanker}, \S\ref{sec:afv} \\
    OE-2 & Telemetría AMQP                   & \S\ref{sec:telemetria} \\
    OE-3 & Contenerización Docker            & \S\ref{sec:docker} \\
    OE-4 & Despliegue Kubernetes             & \S\ref{sec:kubernetes} \\
    OE-5 & Validación escenarios             & Cap.~\ref{cap:resultados} \\
    OE-6 & Caso de uso CyberNEMO             & \S\ref{sec:cybernemo} \\
    \hline
  \end{tabular}
\end{table}
